% GENERATED FILE. Edit canonical lessons, then run npm run generate:books.
% canonical-source-hash: fnv1a64:432ecc972f719aac
% canonical-lessons: GU-C253-shokh, GU-C253-navrash, GU-C253-manvu-enjoy, GU-C253-kantalvu, GU-C253-ajmavvu

\chapter{Hobbies and spare hours}
\label{ch:gu-hobbies-and-spare-hours}
\hlchaptermodality{253}

\begin{chapteropening}
\textbf{By the end of this chapter:} \emph{I can say what hobbies I enjoyed, what bored me and what I tried.}

Complete the last lesson of chapter 253: I can say what hobbies I enjoyed, what bored me and what I tried.
\end{chapteropening}

\section[\texorpdfstring{\gu{શોખ} (śokh)}{શોખ (śokh)}]{\gu{શોખ} (śokh) — hobby, keen interest}

\label{lesson:GU-C253-shokh}



\begin{quote}
Before the new one: say the Gujarati for uniform, then the Gujarati for to retire.
\end{quote}

\subsection*{You'll want to know: \gu{શોખ}}
\textbf{\gu{શોખ}} — \emph{śokh} — \textquotedblleft{}hobby, keen interest\textquotedblright{}.

\begin{quote}
\textbf{\gu{મારા} \gu{દાદાને} \gu{ગાવાનો} \gu{શોખ} \gu{હતો}.} — \emph{mārā dādāne gāvāno śokh hato} — \textquotedblleft{}My grandfather was fond of singing.\textquotedblright{}
\end{quote}

\begin{cousinweb}
\textbf{\gu{શોખ}} (\emph{śokh}) is a loan from Arabic \emph{shawq}, \textquotedblleft{}longing, desire\textquotedblright{}, which reached Gujarati through Persian. Its Arabic root is \emph{sh-w-q} (Lane).
\end{cousinweb}

\begin{grammarlens}[title={What you've built}]
One more word for everyday life at A2.
\end{grammarlens}

\subsection*{Your turn}
\begin{itemize}
\raggedright
  \item \emph{Say it:} \emph{śokh}
  \item \emph{Say it:} \emph{śokh}, once more
  \item \emph{Say it:} \emph{mārā dādāne gāvāno śokh hato} — My grandfather was fond of singing.
  \item \emph{Recall:} say the Gujarati for uniform, then the Gujarati for to retire, then say \emph{śokh} again
\end{itemize}

\begin{tcolorbox}[breakable,colback=teal!4,colframe=teal!35!black,title={Before you move on}]
What does \gu{શોખ} mean? (\textquotedblleft{}Hobby, keen interest\textquotedblright{}.) Say it once more.
\end{tcolorbox}
\section[\texorpdfstring{\gu{નવરાશ} (navrāś)}{નવરાશ (navrāś)}]{\gu{નવરાશ} (navrāś) — leisure, spare time}

\label{lesson:GU-C253-navrash}



\begin{quote}
Before the new one: say the Gujarati for to retire, then the Gujarati for hobby.
\end{quote}

\subsection*{You'll want to know: \gu{નવરાશ}}
\textbf{\gu{નવરાશ}} — \emph{navrāś} — \textquotedblleft{}leisure, spare time\textquotedblright{}.

\begin{quote}
\textbf{\gu{રવિવારે} \gu{મને} \gu{નવરાશ} \gu{હતી}.} — \emph{ravivāre mane navrāś hatī} — \textquotedblleft{}I had spare time on Sunday.\textquotedblright{}
\end{quote}

\begin{cousinweb}
The origin of \textbf{\gu{નવરાશ}} (\emph{navrāś}) is not settled. CDIAL has no entry for it, and our other sources do not trace it.
\end{cousinweb}

\begin{grammarlens}[title={What you've built}]
One more word for everyday life at A2.
\end{grammarlens}

\subsection*{Your turn}
\begin{itemize}
\raggedright
  \item \emph{Say it:} \emph{navrāś}
  \item \emph{Say it:} \emph{navrāś}, once more
  \item \emph{Say it:} \emph{ravivāre mane navrāś hatī} — I had spare time on Sunday.
  \item \emph{Recall:} say the Gujarati for to retire, then the Gujarati for hobby, then say \emph{navrāś} again
\end{itemize}

\begin{tcolorbox}[breakable,colback=teal!4,colframe=teal!35!black,title={Before you move on}]
What does \gu{નવરાશ} mean? (\textquotedblleft{}Leisure, spare time\textquotedblright{}.) Say it once more.
\end{tcolorbox}
\section[\texorpdfstring{\gu{માણવું} (māṇvũ)}{માણવું (māṇvũ)}]{\gu{માણવું} (māṇvũ) — to enjoy, to relish}

\label{lesson:GU-C253-manvu-enjoy}



\begin{quote}
Before the new one: say the Gujarati for hobby, then the Gujarati for leisure.
\end{quote}

\subsection*{You'll want to know: \gu{માણવું}}
\textbf{\gu{માણવું}} — \emph{māṇvũ} — \textquotedblleft{}to enjoy, to relish\textquotedblright{}.

\begin{quote}
\textbf{\gu{અમે} \gu{ગઈકાલે} \gu{રજા} \gu{માણી}.} — \emph{ame gaīkāle rajā māṇī} — \textquotedblleft{}We enjoyed the holiday yesterday.\textquotedblright{}
\end{quote}

Not the taught \textbf{\gu{માનવું}}, \emph{mānvũ}, \textquotedblleft{}to believe\textquotedblright{}.

\begin{cousinweb}
\textbf{\gu{માણવું}} (\emph{māṇvũ}) is a tadbhava. CDIAL lists it, glossed \textquotedblleft{}to enjoy\textquotedblright{}, under Sanskrit \emph{mānáyati}, \textquotedblleft{}esteems, honours\textquotedblright{} (CDIAL 10047).
\end{cousinweb}

\begin{grammarlens}[title={What you've built}]
One more word for everyday life at A2.
\end{grammarlens}

\subsection*{Your turn}
\begin{itemize}
\raggedright
  \item \emph{Say it:} \emph{māṇvũ}
  \item \emph{Say it:} \emph{māṇvũ}, once more
  \item \emph{Say it:} \emph{ame gaīkāle rajā māṇī} — We enjoyed the holiday yesterday.
  \item \emph{Recall:} say the Gujarati for hobby, then the Gujarati for leisure, then say \emph{māṇvũ} again
\end{itemize}

\begin{tcolorbox}[breakable,colback=teal!4,colframe=teal!35!black,title={Before you move on}]
What does \gu{માણવું} mean? (\textquotedblleft{}To enjoy, to relish\textquotedblright{}.) Say it once more.
\end{tcolorbox}
\section[\texorpdfstring{\gu{કંટાળવું} (kanṭāḷvũ)}{કંટાળવું (kanṭāḷvũ)}]{\gu{કંટાળવું} (kanṭāḷvũ) — to get bored, to be fed up}

\label{lesson:GU-C253-kantalvu}



\begin{quote}
Before the new one: say the Gujarati for leisure, then the Gujarati for to enjoy.
\end{quote}

\subsection*{You'll want to know: \gu{કંટાળવું}}
\textbf{\gu{કંટાળવું}} — \emph{kanṭāḷvũ} — \textquotedblleft{}to get bored, to be fed up\textquotedblright{}.

\begin{quote}
\textbf{\gu{ગઈકાલે} \gu{હું} \gu{ઘરે} \gu{કંટાળ્યો}.} — \emph{gaīkāle hũ ghare kanṭāḷyo} — \textquotedblleft{}Yesterday I got bored at home.\textquotedblright{}
\end{quote}

\begin{cousinweb}
The origin of \textbf{\gu{કંટાળવું}} (\emph{kanṭāḷvũ}) is not settled. CDIAL lists no Gujarati verb of this form and meaning, and our other sources do not trace it.
\end{cousinweb}

\begin{grammarlens}[title={What you've built}]
One more word for everyday life at A2.
\end{grammarlens}

\subsection*{Your turn}
\begin{itemize}
\raggedright
  \item \emph{Say it:} \emph{kanṭāḷvũ}
  \item \emph{Say it:} \emph{kanṭāḷvũ}, once more
  \item \emph{Say it:} \emph{gaīkāle hũ ghare kanṭāḷyo} — Yesterday I got bored at home.
  \item \emph{Recall:} say the Gujarati for leisure, then the Gujarati for to enjoy, then say \emph{kanṭāḷvũ} again
\end{itemize}

\begin{tcolorbox}[breakable,colback=teal!4,colframe=teal!35!black,title={Before you move on}]
What does \gu{કંટાળવું} mean? (\textquotedblleft{}To get bored, to be fed up\textquotedblright{}.) Say it once more.
\end{tcolorbox}
\section[\texorpdfstring{\gu{અજમાવવું} (ajmāvvũ)}{અજમાવવું (ajmāvvũ)}]{\gu{અજમાવવું} (ajmāvvũ) — to try out, to test}

\label{lesson:GU-C253-ajmavvu}



\begin{quote}
Before the new one: say the Gujarati for to enjoy, then the Gujarati for to get bored.
\end{quote}

\subsection*{You'll want to know: \gu{અજમાવવું}}
\textbf{\gu{અજમાવવું}} — \emph{ajmāvvũ} — \textquotedblleft{}to try out, to test\textquotedblright{}.

\begin{quote}
\textbf{\gu{મેં} \gu{નવો} \gu{શોખ} \gu{અજમાવ્યો}.} — \emph{mẽ navo śokh ajmāvyo} — \textquotedblleft{}I tried out a new hobby.\textquotedblright{}
\end{quote}

\begin{cousinweb}
\textbf{\gu{અજમાવવું}} (\emph{ajmāvvũ}) is a loan, usually traced to Persian \emph{āzmā-}, the stem of \emph{āzmūdan}, \textquotedblleft{}to test, to try\textquotedblright{}, made into a Gujarati verb with \emph{-āv-}. CDIAL has no entry for it.
\end{cousinweb}

\begin{grammarlens}[title={What you've built}]
One more word for everyday life at A2.
\end{grammarlens}

\subsection*{Your turn}
\begin{itemize}
\raggedright
  \item \emph{Say it:} \emph{ajmāvvũ}
  \item \emph{Say it:} \emph{ajmāvvũ}, once more
  \item \emph{Say it:} \emph{mẽ navo śokh ajmāvyo} — I tried out a new hobby.
  \item \emph{Recall:} say the Gujarati for to enjoy, then the Gujarati for to get bored, then say \emph{ajmāvvũ} again
\end{itemize}

\begin{tcolorbox}[breakable,colback=teal!4,colframe=teal!35!black,title={Before you move on}]
What does \gu{અજમાવવું} mean? (\textquotedblleft{}To try out, to test\textquotedblright{}.) Say it once more.
\end{tcolorbox}
